\newpage
\section{Wstęp}
\subsection{Kontekst i~motywacja}

Zdecentralizowane giełdy (ang. \emph{decentralized exchange}, DEX) umożliwiają wymianę aktywów
cyfrowych za pomocą smart kontraktów. Znacząca część z~tych giełd korzysta z~algorytmu AMM
(ang. \emph{automated market maker}). Ponieważ różne rynki mają niezależne pule płynności,
istnieje możliwość występowania rozbieżności cenowych pomiędzy tymi samymi aktywami.

W~takim przypadku pojawia się możliwość arbitrażu: zakupu aktywu na tańszym rynku i~sprzedaży na
droższym. Dzięki wykorzystaniu tego zjawiska rynki stają się bardziej efektywne: znikają
rozbieżności cenowe. Sama różnica cenowa nie oznacza jednak, że taka okazja jest wykonalna.
Wpływają na nią takie czynniki jak koszty wykonania transakcji, płynność czy kolejność wykonania
transakcji, która w~warunkach konkurencji może ulegać zmianie poprzez mechanizmy znane jako MEV
(ang. \emph{miner extractable value}). W~konsekwencji samo wykrycie rozbieżności cenowej nie
jest wystarczające do oceny możliwości jej wykorzystania w~praktyce. Konieczna jest analiza
czynników wpływających na wykonalność takich okazji, co może stanowić podstawę opracowania
narzędzi analitycznych zabezpieczających potencjalnych uczestników rynku (np. animatorów
rynku), tak aby ich oferty nie były narażone na straty.

\subsection{Cel i~zakres pracy}

Celem projektu jest opracowanie i~implementacja systemu zbierającego i~analizującego
historyczne dane ze zdecentralizowanych giełd w~celu wykrywania rozbieżności cenowych oraz
oceny wykonalności takich okazji. Ma on pełnić rolę narzędzia analitycznego zabezpieczającego
uczestników rynku przy realizacji strategii cenowych. W~ramach projektu opracowano
heurystyczny algorytm oceny wykonalności wykrytych okazji arbitrażowych.

Heurystyczna ocena wykonalności stanowi punkt wyjścia: w~pracy zestawiono prostą heurystykę
progową z~jej wersją skalibrowaną na danych historycznych oraz z~dwoma modelami opartymi na
logice rozmytej, systemem wnioskowania Mamdaniego, którego reguły zapisują wiedzę ekspercką
o~czynnikach wykonalności, oraz siecią neuro-rozmytą ANFIS, w~której parametry reguł są uczone
z~danych. Porównanie tych modeli na wspólnym protokole ewaluacji, z~etykietami uzyskanymi przez
retrospektywną weryfikację on-chain, pozwala ocenić, w~jakim stopniu ocena wykonalności może
być trafniejsza od samego wykrycia rozbieżności cenowej.

Realizacja celu obejmuje następujące zadania:
\begin{itemize}[nosep]
  \item zaprojektowanie i~implementację systemu gromadzącego historyczne stany puli Uniswap~V2
        i~Sushiswap z~węzła Ethereum oraz wykrywającego rozbieżności cenowe,
  \item zbudowanie procedury weryfikacji retrospektywnej, ustalającej, czy rozbieżność została
        skonsumowana on-chain,
  \item opracowanie heurystycznego algorytmu oceny wykonalności i~porównanie go z~modelami
        rozmytymi (system Mamdaniego, ANFIS) na wspólnym protokole ewaluacji,
  \item udostępnienie wyników w~interaktywnym panelu, w~tym w~trybie analizy bieżących danych.
\end{itemize}

\subsection{Założenia techniczne}

Przedmiotem analizy jest sieć Ethereum (mainnet) oraz dwie giełdy działające w~modelu AMM
o~stałym iloczynie: Uniswap~V2 i~Sushiswap. Obie mają identyczne kontrakty puli, tę samą
prowizję i~te same zdarzenia, więc rozbieżność cenową między nimi można liczyć wprost
z~rezerw w~tym samym bloku. Badanie objęło cztery pary tokenów (WETH/USDC, WETH/USDT,
WETH/DAI, WBTC/WETH) w~czterech dwutygodniowych oknach historycznych z~lat 2021--2022,
dobranych tak, aby obejmowały okresy gwałtownych ruchów cen (krach z~maja 2021~r., szczyt
z~listopada 2021~r., upadek Terra/Luna w~maju 2022~r., upadek FTX w~listopadzie 2022~r.).
Okno majowe 2021~r. poprzedza zmianę EIP-1559, pozostałe są po niej, co pozwala sprawdzić
modele w~dwóch reżimach opłat.

Dane pochodzą wyłącznie z~archiwalnego węzła Ethereum przez interfejs JSON-RPC: logi zdarzeń
kontraktów puli (\texttt{eth\_getLogs}), metadane bloków i~paragony transakcji. We wniosku
o~zatwierdzenie tematu przewidziano indeksowanie danych subgraphami usługi The Graph
i~udostępnianie ich backendowi w~postaci zapytań GraphQL. W~trakcie projektowania odstąpiono od
tego rozwiązania z~trzech powodów. Po pierwsze, publiczne subgraphy Uniswap~V2 i~Sushiswap
agregują stan puli w~przedziałach godzinowych i~dziennych, a~ocena wykonalności wymaga rezerw
po każdym bloku. Po drugie, weryfikacja retrospektywna okazji potrzebuje paragonów transakcji
konsumujących (logi \texttt{Transfer}, zużyty gaz), których subgraph nie udostępnia, więc
dostęp do węzła RPC był niezbędny niezależnie od źródła stanów. Po trzecie, logi zdarzeń są
źródłem pierwotnym, z~którego subgraph sam jest budowany; pobranie ich bezpośrednio usuwa
zależność od usługi zewnętrznej i~upraszcza reprodukcję wyników. Pozostałe założenia
z~wniosku zachowano: analizowane są zdarzenia smart kontraktów giełd AMM w~sieci zgodnej
z~EVM, system napisano w~języku TypeScript, a~dane przechowuje relacyjna baza PostgreSQL.

Przyjęto, że transakcja arbitrażowa to atomowa wymiana w~dwóch pulach tej samej pary w~jednej
transakcji (dwa swapy i~transfer, ok. 220\,000 jednostek gazu), a~okazją jest blok, w~którym
rozbieżność przekracza sumę prowizji obu wymian. Arbitraż trójkątny, wielohopowy i~między
sieciami pozostaje poza zakresem. System jest aplikacją lokalną, bez uwierzytelniania
i~bez obsługi reorganizacji łańcucha; analizowane bloki są dawno finalne.

\subsection{Zakres prac własnych a~elementy wykorzystane}

W~ramach projektu zaprojektowano i~zaimplementowano od podstaw: schemat bazy danych
i~migracje; warstwę pobierania danych z~węzła RPC (porcjowanie, wznawianie, rotacja węzłów);
odtwarzanie stanu puli po każdym bloku i~arytmetykę AMM na liczbach całkowitych, w~tym wzór
na optymalną wielkość transakcji; obliczanie cech i~wykrywanie okazji; procedurę weryfikacji
retrospektywnej z~klasyfikacją statusu, beneficjenta i~trasy; heurystykę progową oraz jej
wersję kalibrowaną; implementację ANFIS z~uczeniem gradientowym; protokół ewaluacji
z~podziałem grupowym, bootstrapem i~metrykami; kolejkę zadań, serwer API i~interfejs
użytkownika z~trybem na żywo; zestaw 953 testów automatycznych, potok CI i~skrypty
reprodukcji.

Zdefiniowano cechy $S$, $G$, $L$, $M$
oraz bazę 16 reguł i~parametry termów systemu Mamdaniego, który w~tej pracy zaimplementowano jako moduł oceny i~poddano ewaluacji na etykietach z~weryfikacji. Do wnioskowania
rozmytego użyto biblioteki \texttt{@thi.ng/fuzzy}, do warstwy danych Drizzle ORM i~PostgreSQL,
do API Fastify, do walidacji kontraktów zod, do interfejsu React, Vite, TanStack Query
i~Recharts, a~do testów Vitest. Dekodowanie logów kontraktów oparto na bibliotece ethers.
Wzory arytmetyki puli odtworzono z~kontraktów Uniswap~V2. Narzędzia generatywnej sztucznej inteligencji użyte przy
realizacji projektu opisano w~rozdziale~6.

\subsection{Struktura pracy i~konwencja edytorska}

Rozdział~2 wprowadza pojęcia z~zakresu Ethereum, giełd AMM, arbitrażu i~MEV, logiki rozmytej
oraz oceny klasyfikatorów, a~kończy się przeglądem istniejących rozwiązań i~uzasadnieniem
własnego podejścia. Rozdział~3 opisuje projekt i~implementację systemu: przypadki użycia,
architekturę, model danych, pozyskiwanie danych, wykrywanie okazji, cztery modele oceny,
weryfikację retrospektywną, kolejkę zadań, API i~interfejs. Rozdział~4 przedstawia zbiór
danych, protokół ewaluacji, sprawdzenie poprawności implementacji, porównanie modeli, ich
stabilność oraz wstępną ocenę ekonomiczną, a~zamyka go dyskusja wyników i~ograniczeń.
Rozdział~5 podsumowuje realizację celów i~wskazuje kierunki dalszych prac, rozdział~6
dokumentuje użycie narzędzi GenAI. Załączniki zawierają instrukcję uruchomienia, słownik
pojęć, pełną bazę reguł, opis repozytorium i~listę punktów końcowych API.

W~tekście nazwy plików, tabel, kolumn, zmiennych i~poleceń zapisano czcionką o~stałej
szerokości (np. \texttt{block\_states}, \texttt{eth\_getLogs}), terminy angielskie kursywą
z~podaniem polskiego odpowiednika przy pierwszym wystąpieniu, a~fragmenty kodu źródłowego
TypeScript i~polecenia powłoki jako numerowane listingi (,,Fragment kodu''). Wzory są
numerowane w~nawiasach okrągłych, a~odwołania do literatury podano w~nawiasach kwadratowych
w~kolejności cytowania. Liczby w~tekście zapisano z~przecinkiem dziesiętnym i~odstępem jako
separatorem tysięcy; kwoty w~USD podano według kursu ETH/USD z~bloku, którego dotyczą.
